% WCAG 2.2 Competition Report
\documentclass[12pt,a4paper]{article}

\usepackage[a4paper,margin=2.4cm]{geometry}
\usepackage{fontspec}
\usepackage{polyglossia}
\usepackage{setspace}
\usepackage{array}
\usepackage{booktabs}
\usepackage{longtable}
\usepackage{xcolor}
\usepackage{hyperref}
\usepackage{fancyhdr}

\setdefaultlanguage{vietnamese}
\setotherlanguage{english}

\IfFontExistsTF{Times New Roman}{
  \setmainfont{Times New Roman}
  \setsansfont{Times New Roman}
  \setmonofont{Courier New}
}{
  \setmainfont{Times}
  \setsansfont{Arial}
  \setmonofont{Courier}
}

\definecolor{Primary}{HTML}{111111}
\definecolor{Accent}{HTML}{FCD000}
\definecolor{Muted}{HTML}{666666}

\hypersetup{
  colorlinks=true,
  linkcolor=Primary,
  urlcolor=blue!70!black,
  pdftitle={Báo cáo dự án WCAG 2.2},
  pdfauthor={Nhóm thực hiện},
  pdfsubject={Website/ứng dụng số đạt chuẩn tiếp cận WCAG 2.2}
}

\pagestyle{fancy}
\fancyhf{}
\fancyhead[L]{\small Báo cáo WCAG 2.2}
\fancyhead[R]{\small \thepage}
\renewcommand{\headrulewidth}{0.4pt}
\setlength{\headheight}{14pt}

\setlength{\parindent}{0pt}
\setlength{\parskip}{0.55em}
\setstretch{1.15}
\renewcommand{\labelitemi}{\textbullet}
\renewcommand{\arraystretch}{1.12}

\newcommand{\reporttitle}{Báo cáo sản phẩm dự thi WCAG 2.2}
\newcommand{\projectname}{AccessJobs}
\newcommand{\subtitletext}{Website hỗ trợ tìm việc và tuyển dụng dành cho người khuyết tật}
\newcommand{\reportdate}{17/06/2026}

\begin{document}

\begin{titlepage}
  \centering
  \vspace*{1.4cm}

  {\Large \textbf{CUỘC THI XÂY DỰNG WEBSITE/ỨNG DỤNG SỐ ĐẠT TIÊU CHUẨN TIẾP CẬN WCAG 2.2}\par}
  \vspace{1.8cm}

  {\Huge \bfseries \reporttitle\par}
  \vspace{0.6cm}
  {\LARGE \projectname\par}
  \vspace{0.4cm}
  {\large \subtitletext\par}

  \vfill

  \begin{tabular}{>{\bfseries}p{4.8cm}p{8.1cm}}
    Tên sản phẩm & \projectname \\
    Loại sản phẩm & Website/ứng dụng số tiếp cận \\
    Công nghệ chính & Next.js, TypeScript, Tailwind CSS, next-intl \\
    Tiêu chuẩn mục tiêu & WCAG 2.2 Level AA \\
    Ngày báo cáo & \reportdate \\
  \end{tabular}

  \vfill

  {\large Tài liệu tham khảo chính: \url{https://www.w3.org/TR/WCAG22/}}\par
  \vspace{0.2cm}
  {\large Bộ sơ đồ use case Mermaid được đặt trong thư mục \texttt{usecases/}.}\par
\end{titlepage}

\tableofcontents
\newpage

\section{Tổng quan}
\projectname{} là một nền tảng việc làm và tuyển dụng tập trung vào khả năng tiếp cận cho người khuyết tật. Phiên bản hiện tại được xây dựng theo hướng tĩnh để cố định cấu trúc thông tin, kiểm tra thứ bậc nội dung, hành vi bàn phím, reflow khi phóng to và tính nhất quán của các thành phần giao diện.

Mục tiêu của sản phẩm không chỉ là “có giao diện đẹp”, mà là tạo ra một luồng tìm việc và tuyển dụng rõ ràng, dễ đọc, dễ tab, dễ hiểu và đủ minh chứng để trình bày trước ban giám khảo cuộc thi WCAG 2.2.

\subsection{Phạm vi đánh giá}
Tài liệu này tập trung vào các màn hình và luồng chính đã được triển khai trong sản phẩm:
\begin{itemize}
  \item Trang chủ và trang danh sách việc làm.
  \item Trang chi tiết việc làm, bao gồm mô tả, yêu cầu, quyền lợi và việc làm tương tự.
  \item Đăng nhập, đăng ký, quên mật khẩu và đặt lại mật khẩu.
  \item Hồ sơ ứng viên, CV, việc làm đã ứng tuyển và theo dõi trạng thái ứng tuyển.
  \item Khu vực nhà tuyển dụng: dashboard, tạo tin tuyển dụng và hộp thư CV ứng viên.
\end{itemize}

\subsection{Phương pháp đánh giá}
Nhóm sử dụng ba cách kiểm tra chính:
\begin{itemize}
  \item Kiểm tra thủ công bằng bàn phím: Tab, Shift+Tab, Enter, Space và Esc.
  \item Kiểm tra hiển thị ở nhiều độ rộng màn hình và khi phóng to giao diện.
  \item Rà soát cấu trúc ngữ nghĩa, nhãn biểu mẫu, tiêu đề, liên kết và trạng thái điều hướng.
\end{itemize}

\subsection{Mô hình tác nhân}
\begin{longtable}{>{\bfseries}p{3.5cm}p{4.8cm}p{6.0cm}}
\toprule
Tác nhân & Vai trò chính & Nhiệm vụ nổi bật trong hệ thống \\
\midrule
\endhead
Khách truy cập & Khám phá sản phẩm và xem việc làm công khai & Tìm kiếm, lọc, xem chi tiết việc làm, đăng ký, đăng nhập \\
Người tìm việc & Quản lý hồ sơ cá nhân và ứng tuyển & Cập nhật profile, xem CV, ứng tuyển, theo dõi trạng thái hồ sơ \\
Nhà tuyển dụng & Quản lý đăng tin và ứng viên & Tạo tin tuyển dụng, xem dashboard, kiểm tra CV và cập nhật trạng thái ứng viên \\
\bottomrule
\end{longtable}

\section{Kết quả tổng hợp}
\subsection{Tổng quan trạng thái}
Đánh giá nội bộ hiện tại cho thấy phần lớn các trải nghiệm cốt lõi đã hoàn thiện, còn lại một số mục cần chốt để báo cáo và minh chứng đầy đủ hơn.

\begin{longtable}{>{\bfseries}p{4.2cm}p{2.4cm}p{8.0cm}}
\toprule
Nhóm kết quả & Số lượng & Diễn giải \\
\midrule
\endhead
Đã hoàn thành & 19/28 & Các luồng chính đã có cấu trúc tốt, bàn phím hoạt động ổn, header và navigation đã được chuẩn hoá \\
Hoàn thành một phần & 6/28 & Một số tiêu chí cần bổ sung bằng chứng, tối ưu target size hoặc rà thêm ở các tình huống zoom cao \\
Chưa hoàn thành & 3/28 & Chủ yếu nằm ở bước hoàn thiện hồ sơ minh chứng, ảnh chụp và rà soát cuối \\
\bottomrule
\end{longtable}

\subsection{Điểm mạnh nổi bật}
\begin{itemize}
  \item Luồng tìm việc và luồng tuyển dụng đã được tách rõ theo vai trò.
  \item Các màn hình chính đều có tiêu đề, nội dung, nút hành động và thứ bậc thông tin rõ ràng.
  \item Giao diện đã chú ý nhiều đến reflow, focus visible, target size và mobile responsiveness.
  \item Bộ điều hướng có hỗ trợ role-based navigation, skip link và breadcrumb.
\end{itemize}

\subsection{Các điểm cần lưu ý}
\begin{itemize}
  \item Một số bằng chứng ảnh chụp cần được chèn vào phụ lục trước khi nộp.
  \item Cần rà lại vài control phụ để thống nhất target size và spacing ở zoom cao.
  \item Nên có thêm ghi chú kiểm thử cuối cùng về keyboard, zoom và screen reader.
\end{itemize}

\section{Các điểm đã hoàn thành}
\begin{longtable}{>{\bfseries}p{4.1cm}p{3.1cm}p{8.0cm}}
\toprule
Hạng mục & Luồng/màn hình & Mô tả hoàn thành \\
\midrule
\endhead
Điều hướng công khai & Trang chủ, danh sách việc làm & Người dùng có thể xem nội dung công khai, lọc và đi đến trang chi tiết mà không bị kẹt điều hướng \\
Chi tiết việc làm & \texttt{/jobs/[slug]} & Trang chi tiết ưu tiên mô tả rõ, nội dung reflow tốt hơn và có các phần hỗ trợ tiếp cận cụ thể \\
Đăng nhập / đăng ký & \texttt{/login}, \texttt{/register} & Luồng theo vai trò được tách rõ cho ứng viên và nhà tuyển dụng, nội dung hiển thị phù hợp từng nhóm \\
Hồ sơ ứng viên & \texttt{/profile}, \texttt{/profile/cv}, \texttt{/profile/applied-jobs} & Ứng viên có thể quản lý profile, CV và theo dõi trạng thái ứng tuyển theo dạng dễ đọc \\
Khu vực nhà tuyển dụng & \texttt{/employer/dashboard}, \texttt{/employer/jobs/create} & Nhà tuyển dụng có dashboard, form đăng tin và đường vào danh sách CV ứng viên \\
Hộp thư CV & \texttt{/employer/cv} & Bảng ứng viên, bộ lọc, phân trang và cập nhật trạng thái đã được tổ chức lại cho nghiệp vụ rõ hơn \\
Tương thích bàn phím & Toàn ứng dụng & FOCUS state, thứ tự tab và các nút hành động chính đã được chuẩn hoá để dễ thao tác bằng bàn phím \\
Responsive / reflow & Các màn hình chính & Các khối nội dung dài và hàng thông tin đã được chỉnh lại để giảm cắt chữ khi zoom lớn \\
\bottomrule
\end{longtable}

\section{Bảng chi tiết 28 tiêu chí WCAG 2.2}
\small
\begin{longtable}{p{1.8cm}p{3.0cm}p{0.9cm}p{1.5cm}p{4.5cm}p{4.0cm}}
\toprule
Mã & Tiêu chí & Mức & Trạng thái & Bằng chứng triển khai & Ghi chú \\
\midrule
\endhead
1.1.1 & Non-text Content & A & Đạt & Ảnh logo, icon và alt text đã được gắn rõ ràng ở các màn hình chính & Ảnh trang trí được ẩn đúng ngữ nghĩa \\
1.3.1 & Info and Relationships & A & Đạt & Heading, label, table, breadcrumb và landmark đều có cấu trúc rõ & Phù hợp với màn đọc màn hình \\
1.3.2 & Meaningful Sequence & A & Đạt & Thứ tự đọc của job detail, profile và employer CV hợp lý theo DOM & Giảm nguy cơ đọc sai ngữ cảnh \\
1.3.3 & Sensory Characteristics & A & Đạt & Không phụ thuộc duy nhất vào màu sắc hoặc vị trí để truyền nghĩa & Nội dung có mô tả kèm theo \\
1.3.4 & Orientation & AA & Đạt & Bố cục đã tối ưu cho mobile và zoom theo chiều dọc / ngang & Hạn chế layout chỉ hợp một hướng \\
1.3.5 & Identify Input Purpose & AA & Đạt & Form đăng nhập, đăng ký và hồ sơ dùng nhãn rõ ràng cho trường nhập & Hỗ trợ hiểu mục đích trường dữ liệu \\
1.4.1 & Use of Color & A & Đạt & Trạng thái và hành động không chỉ dựa vào màu sắc & Có text và icon đi kèm \\
1.4.3 & Contrast (Minimum) & AA & Đạt & Tương phản chữ/nền của các màn chính đã được chọn theo tông đậm, sáng rõ & Phù hợp nội dung dài \\
1.4.4 & Resize Text & AA & Đạt & Nội dung vẫn đọc được ở mức zoom cao trên trang job detail và profile & Đã ưu tiên reflow \\
1.4.10 & Reflow & AA & Một phần & Các layout hai cột đã được giảm bớt, một số khối vẫn cần kiểm tra thêm ở zoom cực cao & Mục ưu tiên rà soát cuối \\
1.4.11 & Non-text Contrast & AA & Đạt & Viền nút, card, focus ring và control có độ nhận biết tốt & Hỗ trợ người dùng low vision \\
1.4.12 & Text Spacing & AA & Đạt & Khoảng cách dòng và block nội dung đã được nới hợp lý để dễ đọc & Không phá bố cục khi tăng spacing \\
1.4.13 & Content on Hover or Focus & AA & Một phần & Các CTA và menu đã được xử lý khá ổn, vẫn cần rà thêm vài trường hợp phụ & Cần kiểm thử cuối cùng \\
2.1.1 & Keyboard & A & Đạt & Có thể điều hướng bằng bàn phím qua các trang chính và form & Luồng nghiệp vụ chính đã support \\
2.1.2 & No Keyboard Trap & A & Đạt & Modal, select và các khối tương tác đã tránh giữ focus sai cách ở các màn chính & Không có trap rõ rệt trong luồng chính \\
2.1.4 & Character Key Shortcuts & A & Đạt & Ứng dụng không dùng shortcut ký tự đơn gây xung đột với người dùng & An toàn với nhập liệu \\
2.4.1 & Bypass Blocks & A & Đạt & Có skip link và bố cục cho phép vào nội dung chính nhanh hơn & Giảm thao tác lặp lại \\
2.4.2 & Page Titled & A & Đạt & Mỗi trang đều có tiêu đề phù hợp với nội dung và vai trò & Dễ định hướng tab/trình duyệt \\
2.4.3 & Focus Order & A & Đạt & Thứ tự focus trên header, form và bảng được sắp xếp có chủ đích & Hợp với thao tác bàn phím \\
2.4.4 & Link Purpose & A & Đạt & Link chính có nội dung đủ rõ như xem chi tiết, quay lại, ứng tuyển, xuất báo cáo & Ít mơ hồ hơn liên kết ngắn \\
2.4.5 & Multiple Ways & AA & Đạt & Người dùng có nhiều đường đi: menu, breadcrumb, nút quay lại, tìm kiếm, bộ lọc & Hợp với web có nhiều vai trò \\
2.4.6 & Headings and Labels & AA & Đạt & Tiêu đề và label đã được chuẩn hoá theo từng section và form & Hỗ trợ quét nhanh nội dung \\
2.4.7 & Focus Visible & AA & Đạt & Focus ring và outline hiện rõ ở các control chính & Dễ nhận biết khi tab \\
2.4.11 & Focus Not Obscured & AA & Đạt & Đã xử lý việc focus bị che trên các phần quan trọng như header và bảng dài & Đặc biệt chú ý khi zoom \\
2.4.12 & Focus Not Obscured (Enhanced) & AAA & Một phần & Một số khối phụ vẫn cần rà thêm ở những viewport hẹp hơn & Không phải mục bắt buộc cho nộp AA \\
2.4.13 & Focus Appearance & AAA & Đạt & Focus state có độ tương phản và kích thước đủ dễ thấy trên giao diện chính & Phù hợp tinh thần WCAG 2.2 \\
2.5.7 & Dragging Movements & AA & Đạt & Luồng chính không phụ thuộc thao tác kéo thả bắt buộc & Tránh gây cản trở cho người dùng motor \\
2.5.8 & Target Size (Minimum) & AA & Một phần & Kích thước nút/link đã được cải thiện, vẫn nên rà thêm vài control nhỏ & Mục cần chốt cuối \\
\bottomrule
\end{longtable}
\normalsize

\section{Các điểm chưa hoàn thành}
\begin{longtable}{>{\bfseries}p{2.4cm}p{3.2cm}p{3.7cm}p{3.0cm}p{2.1cm}}
\toprule
Ưu tiên & Vấn đề & Tác động tới người dùng & Hướng xử lý đề xuất & Mức độ \\
\midrule
\endhead
Cao & Chưa chèn đủ ảnh minh chứng vào báo cáo & Làm báo cáo thiếu trực quan, khó đối chiếu với giao diện thực tế & Bổ sung screenshot của các màn hình chính vào phụ lục & Gần hạn \\
Trung bình & Một số control phụ cần rà lại target size & Người dùng touch hoặc zoom cao có thể bấm khó hơn & Tăng đồng nhất chiều cao và khoảng đệm cho control nhỏ & Sớm \\
Trung bình & Cần thêm kiểm thử thực tế screen reader & Chưa có bằng chứng cuối cùng về cách đọc toàn bộ màn hình & Ghi nhận test NVDA/VoiceOver và cập nhật ghi chú & Sớm \\
\bottomrule
\end{longtable}

\section{Kế hoạch hành động}
\begin{longtable}{>{\bfseries}p{0.6cm}p{4.3cm}p{3.4cm}p{1.8cm}p{1.8cm}p{2.0cm}p{1.5cm}}
\toprule
STT & Hành động & Trang/component liên quan & Sprint & Deadline & Phụ trách & Trạng thái \\
\midrule
\endhead
1 & Chèn screenshot cho các màn hình chính vào phụ lục & Report / appendix & S1 & Trước nộp & Nội dung & Chưa làm \\
2 & Rà lại tiêu chí target size cho các nút phụ & Header, card, select & S1 & Trước nộp & UI & Đang làm \\
3 & Kiểm tra lại reflow ở zoom cao trên job detail và profile & \texttt{/jobs/[slug]}, \texttt{/profile} & S1 & Trước nộp & UI + QA & Đang làm \\
4 & Bổ sung ghi chú test bàn phím và focus order & Toàn ứng dụng & S1 & Trước nộp & QA & Đang làm \\
5 & Ghi nhận kiểm thử screen reader tối thiểu & Trang chính, form & S2 & Sau rà soát & QA & Chưa làm \\
6 & Chuẩn hóa bảng WCAG 28 tiêu chí thành bản cuối & Report & S2 & Sau rà soát & Nội dung & Đang làm \\
7 & Đóng gói PDF, link nguồn và phụ lục nộp bài & Report package & S2 & Trước nộp & Nhóm dự án & Chưa làm \\
\bottomrule
\end{longtable}

\section{Kết luận}
\projectname{} đã đạt được nền tảng tốt cho một sản phẩm dự thi WCAG 2.2: luồng nghiệp vụ rõ ràng, giao diện có cấu trúc, hỗ trợ bàn phím tốt và đã xử lý nhiều vấn đề reflow, focus và điều hướng. Điểm mạnh lớn nhất của phiên bản hiện tại là hệ thống đã thể hiện được cách tiếp cận theo vai trò, thay vì chỉ dừng ở một website giới thiệu thông thường.

Trong giai đoạn tiếp theo, nhóm nên tập trung vào hoàn thiện bằng chứng minh họa, rà lại một số control nhỏ và xác nhận thêm bằng kiểm thử thực tế trên screen reader. Sau khi hoàn tất các mục này, báo cáo sẽ đủ chặt để nộp chính thức và dễ thuyết minh trước ban giám khảo.

\vfill
\begin{center}
  \textcolor{Muted}{\small Tài liệu này là bản báo cáo mẫu có thể tiếp tục tinh chỉnh theo minh chứng cuối cùng của đội thi.}
\end{center}

\newpage
\appendix
\section{Phụ lục: Mermaid use case}
Bộ sơ đồ Mermaid được tách thành các file riêng trong thư mục \texttt{docs/wcag-competition-report/usecases/}:
\begin{itemize}
  \item \texttt{usecase-overview.mmd} – sơ đồ tổng quan hệ thống.
  \item \texttt{usecase-guest.mmd} – luồng cho khách truy cập.
  \item \texttt{usecase-candidate.mmd} – luồng cho người tìm việc.
  \item \texttt{usecase-employer.mmd} – luồng cho nhà tuyển dụng.
  \item \texttt{usecase-accessibility.mmd} – luồng hỗ trợ tiếp cận và điều hướng.
\end{itemize}

\section{Phụ lục: Danh sách màn hình đã kiểm thử}
\begin{itemize}
  \item Trang chủ và trang danh sách việc làm.
  \item Trang chi tiết việc làm.
  \item Đăng nhập và đăng ký.
  \item Quên mật khẩu và đặt lại mật khẩu.
  \item Hồ sơ ứng viên và việc làm đã ứng tuyển.
  \item Dashboard nhà tuyển dụng, đăng tin và hộp thư CV.
\end{itemize}

\section{Phụ lục: Gợi ý ảnh minh họa}
Khi chèn ảnh vào bản cuối, ưu tiên dùng ảnh chụp thật của các trang chính và đặt caption ngắn, rõ, đúng ngữ cảnh. Mỗi ảnh nên thể hiện một điểm kiểm thử cụ thể như reflow, keyboard focus, role-based navigation hoặc form accessibility.

\end{document}
